\documentclass[preprint,12pt]{elsarticle}

% ============================================================
% PACKAGES
% ============================================================

\usepackage{amssymb}
\usepackage{amsmath}
\usepackage{amsfonts}
\usepackage{booktabs}
\usepackage{tabularx}
\usepackage{multirow}
\usepackage{makecell}
\usepackage{graphicx}
\usepackage{float}
\usepackage{array}
\usepackage{url}
\usepackage{xcolor}
\usepackage{hyperref}
\usepackage{tikz}
\usetikzlibrary{arrows.meta,positioning,shapes.geometric,calc,backgrounds}

\journal{Computers and Electronics in Agriculture}

% ============================================================
% CUSTOM COMMANDS
% ============================================================

\newcommand{\placeholder}[1]{
\begin{center}
\fbox{
\begin{minipage}[c][4.5cm][c]{0.90\linewidth}
\centering
\textbf{#1}\\[6pt]
\textit{Insert final figure/graph here}
\end{minipage}
}
\end{center}
}

% ============================================================
% DOCUMENT
% ============================================================

\begin{document}

\begin{frontmatter}

% ============================================================
% TITLE
% ============================================================

\title{
Agri-CNX-Edge: A Unified Multi-Task Edge AI Framework for
Agricultural Disease, Pest and Fruit Analysis Using
Heterogeneous Knowledge Distillation
}

% ============================================================
% AUTHORS
% ============================================================

\author[inst1]{Aarya Balwadkar}
\author[inst1]{Afifa Bintul Hasan}
\author[inst1]{Aparna Nair}
\author[inst1]{Eccha Bansal}
\author[inst1]{Dr. Ankur Goyal}

\affiliation[inst1]{
    organization={Department of Computer Science and Engineering,
    Symbiosis Institute of Technology, Symbiosis International (Deemed University)},
    city={Pune},
    state={Maharashtra},
    country={India}
}

% ============================================================
% ABSTRACT
% ============================================================

\begin{abstract}

Agricultural disease diagnosis, pest recognition, fruit classification,
and fruit counting are commonly addressed using separate computer vision
models. Such task-specific systems increase computational redundancy and
complicate deployment on resource-constrained devices. This paper presents
Agri-CNX-Edge, a unified multi-task framework designed to jointly perform
leaf disease classification, pest recognition, fruit classification, and
fruit localization/counting.

The proposed framework combines a pretrained DINOv3-ConvNeXt-Tiny
visual backbone (adapted with LoRA) with hierarchical feature routing,
an FPN-Lite neck, lightweight MBConv-based classification heads, a
decoupled YOLO-style detection head, and ROI-Align-based per-apple
fruit classification.
Specialized teacher models provide task-specific supervisory information
through heterogeneous knowledge distillation: ConvNeXt-Small
classification teachers supply temperature-1 soft probability
distributions (three seeds $\times$ eight test-time augmentation views),
while a YOLO26s teacher supplies bounding-box soft labels with
confidence scores.
Adaptive homoscedastic uncertainty weighting is employed to balance the
heterogeneous task objectives.

The LoRA-adapted configuration reports, on the held-out test set,
100.00\% leaf classification accuracy, 99.56\% pest accuracy, and
95.69\% fruit accuracy, with a counting MAE of 3.53, RMSE of 6.13,
Pearson $r$ of 0.85, and $R^2$ of 0.57 on 400 yield images. The
no-distillation ablation matches or exceeds the proposed configuration
on every metric (fruit test 99.14\%, MAE 3.08), and the frozen backbone
matches it (fruit test 95.69\%, MAE 3.43): under the evaluated
conditions of saturated teachers and ceiling tasks, heterogeneous
distillation does not improve the student, while adaptive task
weighting proves essential (static weighting collapses fruit test
accuracy to 70.26\%). The unified 120.7\,MB model runs fully offline
on a Samsung A17 5G at 5.6--5.7\,s per inference with identical
numerics across providers. All experiments use a single seed ($n{=}1$);
run-to-run variation is estimated from same-method configurations at
${\sim}3$ percentage points on fruit.

A component-level ablation study is organized around removal of knowledge
distillation, static task weighting, LORA/adaptation, and warm-up.
Class-wise test evaluation, counting correlation statistics, spatial
visualization, and single-device on-device benchmarking complete the
evidence. The results establish a working unified multi-task system
with a measured efficiency profile, validate adaptive weighting as
essential, and document the boundary conditions under which
distillation does not help.

\end{abstract}

% ============================================================
% KEYWORDS
% ============================================================

\begin{keyword}

Agricultural artificial intelligence
\sep Multi-task learning
\sep Knowledge distillation
\sep Edge AI
\sep Apple disease detection
\sep Pest recognition
\sep Fruit classification
\sep Fruit localization
\sep Fruit counting
\sep DINOv3-ConvNeXt
\sep YOLO
\sep ROI Align

\end{keyword}

\end{frontmatter}

% ============================================================
% 1. INTRODUCTION
% ============================================================

\section{Introduction}
\label{sec:introduction}

Agriculture plays a critical role in food security, economic development,
and sustainable resource management. Crop diseases, pest infestations,
and inaccurate assessment of fruit production can significantly affect
crop quality and productivity. Computer vision provides an opportunity to
automate agricultural inspection tasks such as disease recognition,
pest identification, fruit classification, fruit localization, and
fruit counting.

Deep learning has become an important approach for agricultural image
analysis because it can learn discriminative visual representations
directly from images. Convolutional neural networks, transfer learning,
object detection, and self-supervised representation learning have
demonstrated strong performance across agricultural applications.

However, many existing systems are designed for a single task. A
disease-classification system, pest-detection system, and fruit-counting
system may therefore require independent models. Deploying multiple
models on resource-constrained hardware results in repeated feature
extraction, higher memory requirements, and increased computational
complexity.

Multi-task learning provides an alternative by allowing several related
tasks to share a common visual representation. Nevertheless, agricultural
tasks have different output structures and loss characteristics. Leaf and
pest recognition are classification problems, while fruit counting is
inherently spatial.

Agri-CNX-Edge addresses this problem using a shared visual backbone,
hierarchical feature routing, an FPN-Lite neck, lightweight
task-specific heads, a decoupled anchor-free YOLO detection head with
ROI-Align-based per-apple classification, heterogeneous knowledge
distillation, and adaptive task weighting.

The principal contributions of this work are:

\begin{itemize}

\item A unified multi-task architecture for leaf disease classification,
pest recognition, fruit classification, and fruit localization/counting.

\item Hierarchical feature routing that assigns fine-grained features
to leaf classification ($f1$), mid-scale features to pest recognition
($f2$), and fused multi-scale features ($P3$ from $f3$+$f4$) to apple
detection, with ROI-Align crops of $P3$ feeding per-apple fruit
classification.

\item A heterogeneous knowledge-distillation strategy that combines
categorical distillation (focal loss on hard labels plus KL divergence
against teacher soft targets at temperature $T=1$) for the three
classification tasks with box-level distillation (CIoU plus objectness
and class alignment against YOLO teacher boxes) for detection.

\item A teacher-supervision pipeline using three classification-teacher
seeds (42, 123, and 456) and eight test-time augmentation views,
producing 24 inference passes per image for the classification soft
labels, complemented by a separately trained YOLO26s apple-detection
teacher supplying bounding-box soft labels with confidence scores.

\item Adaptive homoscedastic uncertainty weighting for balancing
heterogeneous task objectives.

\item A structured ablation framework evaluating the contribution of
knowledge distillation (\texttt{--alpha 1.0 --beta 0.0}), adaptive task
weighting (\texttt{--disable-homo}), LoRA adaptation
(\texttt{--no-lora}), and warm-up (\texttt{--freeze-ep 0}), reported
with the finding that weighting is essential while distillation is
neutral-to-negative under the evaluated conditions.

\item Measured single-device deployment evidence (Samsung A17 5G):
offline latency, peak memory, and model size with identical numerics
across providers.

\end{itemize}

% ============================================================
% 2. RELATED WORK
% ============================================================

\section{Related Work}
\label{sec:related}

\subsection{Agricultural Disease and Pest Recognition}

Deep learning has been extensively investigated for plant disease and
pest recognition. CNN-based architectures and transfer learning methods
have demonstrated strong performance for image-level disease
classification.

Apple disease recognition remains challenging because symptoms vary in
size, shape, severity, and visual appearance. Field images additionally
contain illumination variation, occlusion, complex backgrounds, and
multiple symptoms.

Existing studies have explored architectures such as Xception,
Faster R-CNN, YOLO, and other CNN-based models for agricultural
classification and detection.

\subsection{Fruit Detection and Counting}

Fruit detection and counting are important for orchard monitoring and
yield estimation. Dense fruit clusters, overlapping fruits, foliage
occlusion, and illumination changes make fruit localization challenging.

CenterNet introduced an anchor-free object representation based on
object centers. Such a representation is particularly relevant to fruit
counting because each detected center corresponds to an individual fruit
instance. The implementation in this work follows the same anchor-free
philosophy but uses a decoupled YOLO-style grid head (one prediction per
cell of a $32\times32$ grid) with ROI Align for per-apple classification,
rather than a CenterNet heatmap branch.

YOLO-based agricultural detection systems have also demonstrated the
potential of lightweight object detection for orchard applications.

\subsection{Multi-Task Learning}

Multi-task learning allows several related objectives to share a common
feature representation. This reduces duplicated computation but can
introduce optimization conflicts between tasks.

Different strategies have been proposed for task balancing, including
gradient normalization, gradient conflict mitigation, and uncertainty
based weighting.

In the proposed framework, homoscedastic uncertainty weighting is used
to learn relative task weighting during optimization.

\subsection{Knowledge Distillation}

Knowledge distillation transfers information from a larger teacher model
to a student model. Instead of relying exclusively on hard labels, the
student can learn from the teacher's soft probability distribution.

For spatial tasks, knowledge can also be transferred through intermediate
representations or spatial probability maps.

The proposed framework extends this concept by combining categorical
knowledge distillation for classification tasks with spatial knowledge
distillation for fruit localization.

\subsection{Self-Supervised Visual Representations}

Self-supervised vision models provide transferable representations that
can be adapted to downstream tasks. DINO and DINOv2 demonstrated the
effectiveness of self-supervised learning for obtaining strong visual
features without task-specific supervision.

The proposed architecture uses a DINOv3-ConvNeXt backbone as its shared
visual representation.

% ============================================================
% 3. RESEARCH GAP
% ============================================================

\section{Research Gap}
\label{sec:gap}

Existing agricultural computer vision studies predominantly focus on
individual disease, pest, fruit-detection, or counting tasks. Separate
models require repeated backbone computation and complicate deployment
on resource-constrained systems.

Although multi-task learning can reduce this redundancy, the agricultural
tasks considered in this work have fundamentally different output
structures. Classification requires semantic discrimination, whereas
fruit counting requires spatial localization.

Furthermore, conventional knowledge distillation generally focuses on
classification outputs and does not directly address the transfer of
spatial information for agricultural localization.

Therefore, the research gap addressed in this work is the integration of:

\begin{enumerate}

\item heterogeneous agricultural tasks within a unified architecture,

\item hierarchical feature routing,

\item categorical and spatial knowledge distillation,

\item adaptive task balancing, and

\item an edge-oriented shared computation design.

\end{enumerate}

% ============================================================
% 4. METHODOLOGY
% ============================================================

\section{Methodology}
\label{sec:methodology}

% ------------------------------------------------------------
% 4.1 SYSTEM OVERVIEW
% ------------------------------------------------------------

\subsection{System Overview}

Agri-CNX-Edge consists of five major components:

\begin{enumerate}

\item Dataset curation and preprocessing (cleaning, stratified
70/15/15 splits, task-specific augmentation).

\item Specialized teacher models: ConvNeXt-Small classification teachers
(fine-tuned from a 120-class PDDD checkpoint) and a YOLO26s
apple-detection teacher.

\item Shared DINOv3-ConvNeXt-Tiny student backbone with LoRA adapters.

\item FPN-Lite neck and hierarchical feature routing.

\item Task-specific classification heads, a decoupled YOLO detection
head, and ROI-Align-based per-apple fruit classification.

\end{enumerate}

\begin{figure}[H]
\centering
\resizebox{\textwidth}{!}{\input{paper_figures/tikz_fig01_architecture}}
\caption{Agri-CNX-Edge forward data flow (single idea per diagram).
The shared DINOv3-ConvNeXt-Tiny backbone (stages 0--1 frozen, stages
2--3 LoRA-adapted) routes fine-grained $f1$ to leaf, mid-scale $f2$
to pest, and fused $P3$ to detection, with ROI-Align crops of $P3$
feeding per-apple fruit classification. Dashed arrows are
training-only teacher supervision (losses in Fig.~\ref{fig:loss}).}
\label{fig:arch}
\end{figure}

The input image ($512\times512$) is processed once through the shared
backbone. Fine-grained $f1$ features are routed to leaf classification,
mid-scale $f2$ features to pest recognition, and fused $P3$ features
(from $f3$+$f4$) to apple detection, with ROI-Align crops of $P3$
feeding per-apple fruit classification.

% ------------------------------------------------------------
% 4.2 DATASET
% ------------------------------------------------------------

\subsection{Dataset Curation and Preprocessing}

All classification images pass through an automated cleaning stage:
images smaller than $64$ px or with pixel standard deviation below $5.0$
are removed, and perceptual-hash duplicates (threshold 8) are discarded.
The retained per-class counts are:

\begin{table}[H]
\centering
\caption{Teacher-side dataset cleaning (raw $\rightarrow$ clean).}
\label{tab:dataset}
\begin{tabular}{lcc}
\toprule
\textbf{Class} & \textbf{Raw} & \textbf{Clean} \\
\midrule
Leaf: Apple\_Mosaic & 631 & 622 \\
Leaf: Apple\_\_\_Black\_rot & 621 & 579 \\
Leaf: Alternaria & 631 & 621 \\
Leaf: Healthy & 2276 & 2152 \\
\midrule
Pest: xylotrechus & 319 & 296 \\
Pest: aphids & 314 & 270 \\
Pest: leafhoppers & 573 & 567 \\
Pest: spider\_mite & 299 & 271 \\
\midrule
Fruit: Anthracnose & 295 & 275 \\
Fruit: Black Rot & 542 & 524 \\
Fruit: Healthy & 1424 & 1333 \\
\bottomrule
\end{tabular}
\end{table}

Each classification task is divided by stratified sampling with seed 42
into training, validation, and test partitions in approximately
70/15/15 proportions. The yield (apple detection) data uses its own
train/val split (1093 training and 264 validation images).
Model input resolutions are $224\times224$ for the classification
teachers, $640\times640$ for the YOLO26s teacher, and
$512\times512$ for the student.

\begin{figure}[H]
\centering
\includegraphics[width=\textwidth]{paper_figures/fig02_dataset}
\caption{Dataset curation. Raw versus retained images per class after
automated cleaning (minimum 64~px, pixel standard deviation $\geq 5.0$,
perceptual-hash deduplication threshold 8) for the (a) leaf, (b) pest,
and (c) fruit classification teachers at 224-px input. Stratified
70/15/15 splits (seed 42); yield detection uses its own 1093/264
train/val split at 640-px input.}
\label{fig:dataset}
\end{figure}

% ------------------------------------------------------------
% 4.3 TEACHER ENSEMBLE
% ------------------------------------------------------------

\subsection{Teacher Ensemble}

Specialized teacher models provide task-specific supervision. The three
classification tasks use ConvNeXt-Small teachers fine-tuned from a
120-class PDDD checkpoint (feature dimension 768, dropout 0.5) with
focal loss ($\gamma=2.0$), label smoothing $0.1$, class-balanced
sampling, and minority-class heavy augmentation for classes with fewer
than 400 training images. Training runs for up to 100 epochs with an
8-epoch frozen-backbone phase, a 5-epoch warm-up, and cosine-annealing
warm restarts (backbone LR $1\times10^{-5}$, head LR $1\times10^{-3}$
in phase 1 and $1\times10^{-4}$ in phase 2).

Three independent classification-teacher seeds are used:

\begin{equation}
S=\{42,123,456\}.
\end{equation}

Eight test-time augmentation views (original, flips, and
$\pm15^{\circ}$/$\pm30^{\circ}$ rotations) are evaluated for each
classification-teacher seed. Consequently, the number of teacher
inference passes per image for the classification soft labels is:

\begin{equation}
3\times8=24.
\end{equation}

Per-image soft targets are stored as temperature-1 probabilities
(\texttt{prob\_0}$\ldots$\texttt{prob\_C}). The yield task uses a
separately trained YOLO26s detection teacher (640-pixel input, CIELAB
$a^{*}$ CLAHE enhancement, confidence threshold 0.20, NMS IoU 0.45, up
to 300 detections per image), whose normalized bounding boxes and
confidence scores are stored as \texttt{json\_boxes} soft labels.

% ------------------------------------------------------------
% 4.4 CATEGORICAL KD
% ------------------------------------------------------------

\subsection{Categorical Knowledge Distillation}

For classification tasks, the teacher CSVs already store TTA-averaged
probability vectors at temperature $T=1$, which are used directly with
no further temperature scaling. The student matches them through KL
divergence:

\begin{equation}
L_{KL}
=
D_{KL}
\left(
P_T
\parallel
P_S
\right),
\end{equation}

averaged over samples whose label is not the ignore index ($-1$).
Hard labels are supervised with focal loss ($\gamma=2.0$), and the two
terms are blended with equal weight:

\begin{equation}
L_{cls}
=
\alpha L_{focal}
+
(1-\alpha)L_{KL},
\qquad \alpha=0.5.
\end{equation}

Soft targets contain information regarding the relative confidence of the
teacher among visually similar classes. Images that do not belong to a
given task carry the ignore index and are masked out of that task's
loss.

% ------------------------------------------------------------
% 4.5 SPATIAL KD
% ------------------------------------------------------------

\subsection{Detection Knowledge Distillation}

For apple detection, supervision is box-level rather than heatmap-based.
Each ground-truth or teacher box with center $(c_x,c_y)$ is assigned to
its grid cell:

\begin{equation}
(i,j)=
\left(
\lfloor c_x W \rfloor,
\lfloor c_y H \rfloor
\right),
\qquad W=H=32.
\end{equation}

At positive cells, box regression uses Complete IoU loss while
objectness and single-class scores use binary cross-entropy:

\begin{equation}
L_{det}
=
\lambda_{box}L_{CIoU}
+
\lambda_{obj}L_{obj}
+
\lambda_{cls}L_{cls},
\end{equation}

with:

\begin{equation}
\lambda_{box}=7.5,
\qquad
\lambda_{obj}=1.0,
\qquad
\lambda_{cls}=0.5.
\end{equation}

The loss is computed once against ground-truth boxes and once against
YOLO teacher boxes (coordinates only), with the teacher branch scaled
by:

\begin{equation}
L_{yield}
=
L_{det}^{gt}
+
\beta L_{det}^{teacher},
\qquad \beta=0.5.
\end{equation}

This allows the student to learn spatial fruit-distribution information
from the teacher's predicted boxes and confidence-filtered detections.

% ------------------------------------------------------------
% 4.6 BACKBONE
% ------------------------------------------------------------

\subsection{Student Backbone}

The proposed student uses a DINOv3-ConvNeXt-Tiny backbone as the shared
visual feature extractor, instantiated through \texttt{timm} with
\texttt{features\_only=True}. DINOv3 pretraining weights
(LVD-1689M) are remapped onto the \texttt{timm} naming scheme and
verified for key alignment and forward-pass shapes.

The backbone generates hierarchical feature representations at
different spatial resolutions for a $512\times512$ input:

\begin{equation}
f1:[B,96,128,128],\
f2:[B,192,64,64],\
f3:[B,384,32,32],\
f4:[B,768,16,16].
\end{equation}

Stages 0 and 1 are fully frozen. Stages 2 and 3 carry custom LoRA
adapters on the ConvNeXt MLP pointwise layers (\texttt{mlp.fc1},
\texttt{mlp.fc2}) with rank 8, $\alpha=16$ (scaling 2.0), and dropout
0.05; all other backbone parameters are frozen, leaving only the LoRA
$A$/$B$ matrices trainable. The frozen-backbone configuration instead
applies no LoRA layers at all.

The routing strategy is:

\begin{equation}
f1\rightarrow Leaf,
\qquad
f2\rightarrow Pest,
\end{equation}

while $f3$ and $f4$ are fused by the FPN-Lite neck into $P3$
($[B,128,32,32]$):

\begin{equation}
P3\rightarrow Apple\ Detection,
\qquad
ROIAlign(P3)\rightarrow Fruit\ Classification.
\end{equation}

This design provides fine-grained texture to leaf classification,
mid-scale morphology to pest recognition, and multi-scale fused features
to detection, with per-apple crops for fruit disease classification.

\begin{figure}[H]
\centering
\resizebox{\textwidth}{!}{\input{paper_figures/tikz_fig03_routing}}
\caption{Hierarchical feature routing. Stage shapes and frozen/LoRA
status are as implemented: fine-grained $f1$ serves leaf texture,
mid-scale $f2$ pest morphology, fused $P3$ apple detection
(confidence 0.3, NMS IoU 0.45), and ROI-Align $7\times7$ crops of $P3$
per-apple fruit classification from ground-truth boxes at train time
and decoded boxes at inference.}
\label{fig:routing}
\end{figure}

% ------------------------------------------------------------
% 4.7 CLASSIFICATION HEADS
% ------------------------------------------------------------

\subsection{Lightweight Classification Heads}

The classification heads share an MBConv-based architecture: a
$1\times1$ reduction to a 256-channel bottleneck, one MBConv block
(expansion 4, depthwise convolution, squeeze-and-excitation with ratio
0.25, projection), global average pooling, dropout with $p=0.3$, and a
linear classifier.

Each head consumes its routed feature map:

\begin{equation}
Leaf:f1\ (96ch)\rightarrow 4\ classes,
\quad
Pest:f2\ (192ch)\rightarrow 4\ classes,
\quad
Fruit:ROI\ crops\ (128ch,7\times7)\rightarrow 3\ classes.
\end{equation}

The leaf classes are Apple\_Mosaic, Apple\_\_\_Black\_rot, Alternaria,
and Healthy; the pest classes are xylotrechus, aphids, leafhoppers, and
spider\_mite; the fruit classes are Anthracnose, Black Rot, and
Healthy. Separate classification layers are used for leaf, pest, and
fruit prediction.

% ------------------------------------------------------------
% 4.8 FRUIT LOCALIZATION
% ------------------------------------------------------------

\subsection{YOLO-Style Apple Detection and Per-Apple Classification}

The detection branch follows an anchor-free, single-class, decoupled
YOLO design operating on $P3$ ($[B,128,32,32]$, stride 16). Three
separate branches predict box coordinates, objectness, and class
scores:

\begin{equation}
\hat{D}
\in
\mathbb{R}^{B\times 6\times 32\times 32},
\end{equation}

where the six channels are $[c_x,c_y,w,h,obj,cls]$ and the sigmoid is
applied at loss and post-processing time, not in the head. Objectness
and class biases are initialized from a prior probability of 0.01.

Decoding applies sigmoid activation, keeps cells with
$obj\times cls$ above:

\begin{equation}
conf_{thresh}=0.3,
\end{equation}

and applies non-maximum suppression with:

\begin{equation}
nms_{iou}=0.45.
\end{equation}

For per-apple fruit classification, boxes (ground-truth boxes during
training, decoded boxes at inference) are ROI-Aligned from $P3$ into
$7\times7$ crops with spatial scale $32/512$ and classified by the
fruit head. Per-image apple counts are obtained by counting decoded
boxes, and counting error is reported as MAE/RMSE on the yield
validation set.

% ============================================================
% 5. MULTI-TASK OPTIMIZATION
% ============================================================

\section{Multi-Task Optimization}
\label{sec:optimization}

The four tasks produce heterogeneous loss functions and therefore
require appropriate task balancing.

Homoscedastic uncertainty weighting is employed.

For task $t$, the weighter maintains a learned log-variance:

\begin{equation}
s_t=\log(\sigma_t^2),
\qquad initialized\ to\ 0.
\end{equation}

The multi-task objective is:

\begin{equation}
L_{MTL}
=
\sum_{t\in\mathcal{T}}
\left(
e^{-s_t}L_t
+
\frac{1}{2}s_t
\right),
\end{equation}

where:

\begin{equation}
\mathcal{T}
=
\{
leaf,\ pest,\ fruit,\ yield
\}.
\end{equation}

No hard clamping interval is imposed on $s_t$ in the implementation;
reported terminal weights are precisions $e^{-s_t}$. Knowledge
distillation needs no extra global coefficients because it is already
contained inside the per-task losses ($\alpha$ blending for
classification, $\beta$ scaling of the teacher box branch). Hence the
complete objective is simply:

\begin{equation}
L_{total}
=
L_{MTL}.
\end{equation}

Training is two-phase: Phase~1 freezes the LoRA adapters for 5 epochs
(heads and neck only), and Phase~2 unfreezes LoRA with two AdamW
parameter groups (LoRA at $1\times10^{-5}$, heads/neck/weighter at
$1\times10^{-4}$), cosine annealing, gradient accumulation over 4 steps
(effective batch 32), gradient clipping at 1.0, and early stopping with
patience 20. The static-weight ablation replaces the weighter with a
plain sum of the four task losses.

\begin{figure}[H]
\centering
\resizebox{\textwidth}{!}{\input{paper_figures/tikz_fig04_loss_training}}
\caption{(a) Heterogeneous knowledge distillation for one classification
head: teacher soft targets ($T{=}1$ CSVs) meet student logits in KL
divergence alongside focal loss on hard labels, while detection
distills YOLO teacher boxes through a 0.5-weighted branch.
(b) Two-phase training timeline with best-validation epochs (186 LoRA,
194 frozen), cosine annealing, and ONNX export. Single seed ($n{=}1$);
early-stopping patience 20.}
\label{fig:loss}
\end{figure}

% ============================================================
% 6. EXPERIMENTAL SETUP
% ============================================================

\section{Experimental Setup}
\label{sec:experimental}

\subsection{Implementation Details}

The framework is implemented in PyTorch (timm backbones,
Albumentations augmentation, scikit-learn metrics) with ONNX export
(opset 17) of leaf logits, pest logits, and the raw detection grid;
per-apple fruit classification runs client-side from ROI-Align crops.

Training runs for up to 200 student epochs with batch size 8 and
accumulation over 4 steps (effective batch 32) on CUDA hardware when
available (prior runs used an RTX A4000-class GPU); automatic mixed
precision is disabled because CIoU and KL terms proved unstable in
FP16.

The AdamW optimizer is used with a weight decay of:

\begin{equation}
1\times10^{-4}.
\end{equation}

Two learning-rate groups are used: LoRA adapter parameters at:

\begin{equation}
LR_{lora}=1\times10^{-5},
\end{equation}

and classification heads, FPN-Lite neck, and uncertainty weights at:

\begin{equation}
LR_{heads}=1\times10^{-4}.
\end{equation}

Phase~1 (5 epochs) trains heads and neck with LoRA frozen; Phase~2
unfreezes LoRA. A cosine annealing schedule is used with:

\begin{equation}
\eta_{min}=1\times10^{-7},
\end{equation}

with gradient clipping at norm 1.0 and early stopping (patience 20,
min-delta $1\times10^{-4}$) on total validation loss.

% ------------------------------------------------------------
% METRICS
% ------------------------------------------------------------

\subsection{Evaluation Metrics}

Classification accuracy is calculated as:

\begin{equation}
Accuracy=
\frac{N_{correct}}{N_{total}}.
\end{equation}

Additional class-wise metrics should include precision, recall, and
F1-score.

Precision is:

\begin{equation}
Precision=
\frac{TP}{TP+FP}.
\end{equation}

Recall is:

\begin{equation}
Recall=
\frac{TP}{TP+FN}.
\end{equation}

The F1-score is:

\begin{equation}
F1=
2
\frac{Precision\times Recall}
{Precision+Recall}.
\end{equation}

For fruit counting, Mean Absolute Error is:

\begin{equation}
MAE=
\frac{1}{M}
\sum_{j=1}^{M}
|\hat{N}_j-N_j|.
\end{equation}

Root Mean Squared Error is:

\begin{equation}
RMSE=
\sqrt{
\frac{1}{M}
\sum_{j=1}^{M}
(\hat{N}_j-N_j)^2
}.
\end{equation}

% ============================================================
% 7. RESULTS
% ============================================================

\section{Results}
\label{sec:results}

% ------------------------------------------------------------
% 7.1 FINAL PERFORMANCE
% ------------------------------------------------------------

\subsection{Final-Round Overall Performance}

The best-validation-epoch performance of the LoRA-adapted (proposed)
model is presented in Table~\ref{tab:final_results}. All accuracies are
validation accuracies at epoch 186 (best total validation loss); the
counting errors are computed on the 400-image yield validation set.

\begin{table}[H]
\centering
\caption{Best-validation-epoch performance of the proposed
Agri-CNX-Edge model (LoRA adapters, epoch 186).}
\label{tab:final_results}
\begin{tabular}{lcc}
\toprule
\textbf{Task} & \textbf{Metric} & \textbf{Performance} \\
\midrule
Leaf disease classification & Accuracy (val) & 99.67\% \\
Pest recognition & Accuracy (val) & 100.00\% \\
Fruit classification & Accuracy (val) & 96.30\% \\
Fruit localization/counting & Count MAE & 3.5675 \\
Fruit localization/counting & Count RMSE & 6.3767 \\
\bottomrule
\end{tabular}
\end{table}

\begin{figure}[H]
\centering
\includegraphics[width=\textwidth]{paper_figures/fig05_final_performance}
\caption{Best-validation-epoch snapshot of the proposed LoRA-adapted
model (epoch 186). (a) Leaf, pest, and fruit validation accuracies
against the 100\% reference. (b) Apple counting MAE/RMSE on the
400-image yield validation set (confidence 0.3, NMS IoU 0.45). Single
seed ($n{=}1$).}
\label{fig:final}
\end{figure}

% ------------------------------------------------------------
% 7.2 FROZEN VS PROPOSED
% ------------------------------------------------------------

\subsection{Frozen Versus Proposed Configuration}

The LoRA-adapted (proposed) configuration is compared against the
frozen-backbone configuration (no LoRA layers; best-validation epoch
194). Checkpoint inspection confirms the attribution: the proposed
checkpoint contains trained \texttt{lora\_A}/\texttt{lora\_B} adapters
while the frozen checkpoint contains none.

\begin{table}[H]
\centering
\caption{Comparison between frozen (epoch 194) and proposed (epoch 186)
configurations. Accuracies are best-epoch validation accuracies.}
\label{tab:frozen_proposed}
\begin{tabular}{lccc}
\toprule
\textbf{Metric}
&
\textbf{Frozen}
&
\textbf{Proposed}
&
\textbf{Change}
\\
\midrule
Leaf Accuracy
& 99.50\%
& 99.67\%
& +0.17 pp
\\

Pest Accuracy
& 100.00\%
& 100.00\%
& 0.00 pp
\\

Fruit Accuracy
& 98.42\%
& 96.30\%
& -2.12 pp
\\

Count MAE
& 3.4300
& 3.5675
& +0.1375
\\

Count RMSE
& 6.0017
& 6.3767
& +0.3750
\\

\bottomrule
\end{tabular}
\end{table}

The frozen configuration achieves higher fruit classification accuracy
(98.42\% vs 96.30\%) and lower counting errors (Count MAE 3.4300 vs
3.5675; Count RMSE 6.0017 vs 6.3767) at its best-validation epoch.

Leaf accuracy is 0.17 percentage points higher with LoRA adaptation,
while pest accuracy is unchanged at 100.00\% in both configurations.
The inversion on fruit and counting tasks is reported as observed; the
ablation study is designed to isolate which training component accounts
for it.

\begin{figure}[H]
\centering
\includegraphics[width=\textwidth]{paper_figures/fig06_frozen_vs_proposed}
\caption{Frozen-backbone (best epoch 194, circles) versus proposed
LoRA-adapted (best epoch 186, diamonds) configurations with paired
deltas. LoRA adds $+0.17$~pp on leaf but loses $2.12$~pp on fruit and
adds $+0.14$/$+0.38$ counting error, motivating the ablation study.
Single seed ($n{=}1$).}
\label{fig:compare}
\end{figure}

% ------------------------------------------------------------
% 7.3 CONVERGENCE
% ------------------------------------------------------------

\subsection{Training Convergence}

The proposed (LoRA) configuration achieved its best total validation
loss at epoch 186, while the frozen configuration peaked at epoch 194.

\begin{table}[H]
\centering
\caption{Training convergence summary (best-validation-loss epochs).}
\label{tab:convergence}
\begin{tabular}{lcc}
\toprule
\textbf{Training Statistic} & \textbf{Proposed} & \textbf{Frozen}\\
\midrule
Best epoch & 186 & 194\\
Best validation loss & -1.9251 & -1.9271\\
Training loss at best epoch & -0.7437 & -0.6637\\
\bottomrule
\end{tabular}
\end{table}

\begin{figure}[H]
\centering
\includegraphics[width=\textwidth]{paper_figures/fig07_convergence}
\caption{Training dynamics over 200 epochs (thin raw, thick 7-epoch
smooth). (a) Mean training loss and homoscedastic-weighted total
validation loss with best-validation epochs marked (186 proposed, 194
frozen). (b) Per-task validation accuracies (solid proposed, dashed
frozen); fruit is the noisiest task. Single seed ($n{=}1$).}
\label{fig:conv}
\end{figure}

% ------------------------------------------------------------
% 7.4 TASK WEIGHTS
% ------------------------------------------------------------

\subsection{Learned Task Weights}

The terminal (best-epoch) task precisions $e^{-s_t}$ of the proposed
configuration are presented in Table~\ref{tab:weights_proposed}.

\begin{table}[H]
\centering
\caption{Terminal task precisions of the proposed configuration
(epoch 186).}
\label{tab:weights_proposed}
\begin{tabular}{lc}
\toprule
\textbf{Task} & \textbf{Terminal Weight}\\
\midrule
Fruit & 18.8170\\
Leaf & 6.9591\\
Pest & 5.7344\\
Yield & 0.1780\\
\bottomrule
\end{tabular}
\end{table}

For comparison, the frozen configuration (epoch 194) reports:

\begin{table}[H]
\centering
\caption{Terminal task precisions of the frozen configuration
(epoch 194).}
\label{tab:weights_frozen}
\begin{tabular}{lc}
\toprule
\textbf{Task} & \textbf{Terminal Weight}\\
\midrule
Fruit & 18.4710\\
Leaf & 7.1471\\
Pest & 5.7203\\
Yield & 0.1760\\
\bottomrule
\end{tabular}
\end{table}

These weights represent learned optimization coefficients and should not
be interpreted directly as measures of task accuracy or task importance.

\begin{figure}[H]
\centering
\includegraphics[width=\textwidth]{paper_figures/fig08_task_weights}
\caption{Learned homoscedastic task precisions $e^{-s_t}$ for (a)
proposed and (b) frozen configurations with best-epoch values
annotated. Precisions are optimization coefficients balancing
heterogeneous losses (single seed, $n{=}1$), not measures of task
importance.}
\label{fig:weights}
\end{figure}

% ============================================================
% 8. ABLATION STUDY
% ============================================================

\section{Ablation Study}
\label{sec:ablation}

To quantify the contribution of the proposed components, four ablation
configurations are considered, each launched via explicit run flags:

\begin{enumerate}

\item \textbf{A -- No KD:} knowledge distillation is removed
(\texttt{--alpha 1.0 --beta 0.0}, hard labels only).

\item \textbf{B -- Static Weights:} adaptive task weighting is replaced
with a plain sum of task losses (\texttt{--disable-homo}).

\item \textbf{C -- No LORA:} LoRA adapters are disabled entirely
(\texttt{--no-lora}, empty LoRA stage list).

\item \textbf{D -- No Warm-up:} the warm-up training stage is removed
(\texttt{--freeze-ep 0}, LoRA active from epoch 1).

\end{enumerate}

Ablation histories for A, C, and D cover the full 200-epoch schedule;
ablation B currently has 103 logged epochs and must be resumed or
re-run to completion before its row can be reported. The final results
should be inserted after all four experimental runs are completed;
validation accuracies below are best-epoch validation values, consistent
with the main results.

\begin{table}[H]
\centering
\caption{Component-level ablation study.}
\label{tab:ablation}
\small
\begin{tabular}{lccccc}
\toprule
\textbf{Configuration}
&
\textbf{Leaf Acc.}
&
\textbf{Pest Acc.}
&
\textbf{Fruit Acc.}
&
\textbf{MAE}
&
\textbf{RMSE}
\\
\midrule

Proposed
& 99.67\%
& 100.00\%
& 96.30\%
& 3.5675
& 6.3767
\\

No KD
& 100.00\%
& 100.00\%
& 98.47\%
& 3.0175
& 5.2714
\\

Static Weights
& 99.35\%
& 98.33\%
& 81.43\%
& 3.5825
& 6.0640
\\

No LORA
& 99.67\%
& 99.58\%
& 94.69\%
& 3.6300
& 6.2032
\\

No Warm-up
& 99.67\%
& 100.00\%
& 97.18\%
& 3.6575
& 6.6695
\\

Frozen
& 99.50\%
& 100.00\%
& 98.42\%
& 3.4300
& 6.0017
\\

\bottomrule
\end{tabular}
\end{table}

\begin{figure}[H]
\centering
\includegraphics[width=\textwidth]{paper_figures/fig09_ablation}
\caption{Component-level ablation at best-validation epochs.
(a) Min-max-normalised radar (higher is better; errors inverted) and
(b) exact values: No-KD (A) matches or beats the proposed model on
every metric, while static weighting (B) collapses fruit accuracy.
*B logged only 103/200 epochs and must be completed; all runs
single-seed ($n{=}1$).}
\label{fig:ablation}
\end{figure}

% ============================================================
% 9. CLASS-WISE ANALYSIS
% ============================================================

\section{Class-wise Classification Analysis}
\label{sec:classwise}

\subsection{Leaf Disease Classification}

Held-out test predictions (623 images; proposed configuration evaluated
from the epoch-200 checkpoint as only \texttt{latest\_checkpoint.pth}
was present on the evaluation machine) give the confusion matrix in
Fig.~\ref{fig:leaf_cm} and class-wise scores in
Table~\ref{tab:leaf_classwise}.

\begin{figure}[H]
\centering
\includegraphics[width=\columnwidth]{paper_figures/fig10_leaf_cm}
\caption{Leaf disease confusion matrix on the held-out test set
(623 images; counts). All four classes score 1.00 precision, recall,
and F1. Single seed ($n{=}1$).}
\label{fig:leaf_cm}
\end{figure}

\begin{table}[H]
\centering
\caption{Class-wise leaf disease results on the held-out test set
(supports 101/103/90/329).}
\label{tab:leaf_classwise}
\begin{tabular}{lccc}
\toprule
\textbf{Class} &
\textbf{Precision} &
\textbf{Recall} &
\textbf{F1-score}
\\
\midrule
Apple\_Mosaic & 1.0000 & 1.0000 & 1.0000\\
Apple\_\_\_Black\_rot & 1.0000 & 1.0000 & 1.0000\\
Alternaria & 1.0000 & 1.0000 & 1.0000\\
Healthy & 1.0000 & 1.0000 & 1.0000\\
\bottomrule
\end{tabular}
\end{table}

\subsection{Pest Recognition}

Fig.~\ref{fig:pest_cm} and Table~\ref{tab:pest_classwise} report the
held-out pest test (225 images). A single leafhoppers image is
predicted as xylotrechus; all other images are correct.

\begin{figure}[H]
\centering
\includegraphics[width=\columnwidth]{paper_figures/fig11_pest_cm}
\caption{Pest confusion matrix on the held-out test set (225 images;
counts). The only error is one leafhoppers image predicted as
xylotrechus. Single seed ($n{=}1$).}
\label{fig:pest_cm}
\end{figure}

\begin{table}[H]
\centering
\caption{Class-wise pest results on the held-out test set (supports
52/51/82/40).}
\label{tab:pest_classwise}
\begin{tabular}{lccc}
\toprule
\textbf{Class} &
\textbf{Precision} &
\textbf{Recall} &
\textbf{F1-score}
\\
\midrule
xylotrechus & 0.9811 & 1.0000 & 0.9905\\
aphids & 1.0000 & 1.0000 & 1.0000\\
leafhoppers & 1.0000 & 0.9878 & 0.9939\\
spider\_mite & 1.0000 & 1.0000 & 1.0000\\
\bottomrule
\end{tabular}
\end{table}

\subsection{Fruit Classification}

Fruit is evaluated as a head diagnostic on ground-truth-box ROIs
(232 test ROIs). Fig.~\ref{fig:fruit_cm} and
Table~\ref{tab:fruit_classwise} report the held-out test. No Healthy
fruit ROI is present in validation or test (support 0), so Healthy
generalization is untested and the row is shown as ``---''; the
per-present-class macro-F1 is 0.95.

\begin{figure}[H]
\centering
\includegraphics[width=\columnwidth]{paper_figures/fig12_fruit_cm}
\caption{Fruit confusion matrix on ground-truth-box ROIs (232 test
ROIs; counts). Nine Black Rot ROIs are predicted as Anthracnose and
one as Healthy; the Healthy row has no test ROIs (``---''). Single
seed ($n{=}1$).}
\label{fig:fruit_cm}
\end{figure}

\begin{table}[H]
\centering
\caption{Class-wise fruit results on ground-truth-box ROIs (supports
55/177/0; ``---'' denotes no test ROIs).}
\label{tab:fruit_classwise}
\begin{tabular}{lccc}
\toprule
\textbf{Class} &
\textbf{Precision} &
\textbf{Recall} &
\textbf{F1-score}
\\
\midrule
Anthracnose & 0.8594 & 1.0000 & 0.9244\\
Black Rot & 1.0000 & 0.9435 & 0.9709\\
Healthy & --- & --- & ---\\
\bottomrule
\end{tabular}
\end{table}

% ============================================================
% 10. FRUIT COUNTING
% ============================================================

\section{Fruit Counting and Spatial Evaluation}
\label{sec:counting}

The fruit localization branch is evaluated using counting error metrics
and qualitative spatial predictions.

\subsection{Fruit Counting Results}

The counting evaluation of the proposed configuration on the 400-image
yield set (evaluated from the epoch-200 checkpoint) is reported in
Table~\ref{tab:counting} with the predicted-versus-ground-truth scatter
in Fig.~\ref{fig:counting}. Correlation is good ($r{\approx}0.85$) with
a systematic undercounting bias (mean signed error $-3.1$ apples per
image), consistent with a strict confidence threshold on dense
clusters rather than a ranking failure.

\begin{table}[H]
\centering
\caption{Fruit counting performance (proposed, 400 yield images,
confidence 0.3, NMS IoU 0.45).}
\label{tab:counting}
\begin{tabular}{lc}
\toprule
\textbf{Metric} & \textbf{Value}\\
\midrule
Count MAE & 3.5300\\
Count RMSE & 6.1323\\
Pearson correlation & 0.8508\\
$R^2$ & 0.5700\\
Mean signed error & -3.1000\\
\bottomrule
\end{tabular}
\end{table}

\subsection{Predicted Versus Ground-Truth Counts}

\begin{figure}[H]
\centering
\includegraphics[width=\columnwidth]{paper_figures/fig13_counting}
\caption{Predicted versus ground-truth apple counts on the 400-image
yield set (points; dashed identity). Pearson $r{=}0.85$, $R^2{=}0.57$,
mean signed error $-3.1$ (systematic undercounting). Confidence 0.3,
NMS IoU 0.45; single seed ($n{=}1$).}
\label{fig:counting}
\end{figure}

\subsection{Spatial Localization}

Representative apple detection examples showing ground-truth boxes
versus predicted boxes (confidence threshold 0.3, NMS IoU 0.45) with
per-image ground-truth/predicted counts are shown in
Fig.~\ref{fig:boxes}.

\begin{figure}[H]
\centering
\includegraphics[width=\textwidth]{paper_figures/fig14_boxes}
\caption{Ground-truth (solid) versus predicted (dashed) apple boxes
with per-image counts: (a) sparse case, (b) dense cluster. Confidence
0.3, NMS IoU 0.45.}
\label{fig:boxes}
\end{figure}

% ============================================================
% 11. QUALITATIVE RESULTS
% ============================================================

\section{Qualitative Results}
\label{sec:qualitative}

Representative detection and counting cases, from sparse trees to
dense occluded clusters, are shown in Fig.~\ref{fig:qualitative}.
Single-image leaf, pest, and fruit classification examples are covered
by the confusion-matrix analysis in
Figs.~\ref{fig:leaf_cm}--\ref{fig:fruit_cm}; dedicated per-image
classification montages are deferred to future work.

\begin{figure}[H]
\centering
\includegraphics[width=\textwidth]{paper_figures/fig15_qualitative}
\caption{Qualitative detection and counting cases (a)--(f), ordered
from sparse to dense, with ground-truth/predicted counts. Solid boxes
are ground truth, dashed boxes are predictions (confidence 0.3, NMS
IoU 0.45). Dense clusters show the undercounting bias reported in
Table~\ref{tab:counting}.}
\label{fig:qualitative}
\end{figure}

% ============================================================
% 12. EDGE DEPLOYMENT
% ============================================================

\section{Edge Deployment Analysis}
\label{sec:edge}

The architecture is designed to reduce computational redundancy by
sharing the visual backbone across multiple tasks. The trained student
exports to ONNX (opset 17, \texttt{adc\_student\_full.onnx}, 120.7\,MB)
with outputs for leaf logits, pest logits, and the raw detection grid;
per-apple fruit classification is performed client-side from ROI-Align
crops. The torch checkpoint (\texttt{best\_model.pth}, epoch 186) is
148.5\,MB. The backbone contributes 28.28\,M parameters, of which the
LoRA adapters are 0.46\,M (1.63\%).

On-device evidence was collected on a single representative phone,
Samsung A17 5G (SM-A176B, Exynos 1330, 6\,GB RAM, Android 16, ONNX
Runtime 1.24.3), over adb-over-WiFi in airplane mode at fixed low
brightness with cooldown between modes. The uniform bench-v2 protocol
runs an accuracy-once pass (50 images: 30 yield val with ground-truth
counts plus 20 pest images, one pass each per provider mode) and a
latency-repeat pass (10 green yield images, 10 warmup $+$ 100 measured
inferences per provider mode, batch 1, wall clock covering both ONNX
passes). CPU-only and NNAPI-preferred modes use the same APK, model
file, images, and code on every run.

\begin{table}[H]
\centering
\caption{On-device results on Samsung A17 5G (single-device scope).
Latency is wall clock per inference, batch 1. Accuracy-once uses 30
yield val images plus 20 pest images, one pass each per mode.}
\label{tab:edge}
\begin{tabular}{lc}
\toprule
\textbf{Metric}
&
\textbf{Measured}
\\
\midrule

Backbone parameters
& 28.28\,M (LoRA 0.46\,M)
\\

ONNX file (\texttt{adc\_student\_full.onnx})
& 120.7\,MB
\\

Torch checkpoint (\texttt{best\_model.pth}, epoch 186)
& 148.5\,MB
\\

Leaf teacher file (\texttt{leaf\_best\_teacher.pth})
& 188.8\,MB
\\

Fruit teacher file (\texttt{fruit\_best\_teacher.pth})
& 188.8\,MB
\\

Yield teacher file (\texttt{best.pt}, YOLO26s)
& 19.4\,MB
\\

Pest teacher file
& not present on disk (excluded, not estimated)
\\

Yield counting MAE, phone (30 val, CPU / NNAPI)
& 2.9 / 2.9
\\

Pest agreement, phone vs laptop (20 images)
& 20/20 both modes
\\

CPU vs NNAPI prediction differences (50 images)
& 0 (identical)
\\

Latency mean, CPU / NNAPI ($n{=}100$ measured $+$ 10 warmup)
& 5749 / 5604\,ms
\\

Latency p50, CPU / NNAPI
& 7041 / 6763\,ms
\\

Latency p95, CPU / NNAPI ($n{=}100$, approximate tail)
& 7381 / 7406\,ms
\\

Peak PSS (stable range)
& 765--807\,MB
\\

Battery
& coarse drain only (84\%$\rightarrow$75\% across the v2 tap); no power claim
\\

Wall-power consumption (W)
& not measured (external meter required; not estimated)
\\

Offline operation
& Yes (on-device ONNX Runtime, no server)
\\

\bottomrule
\end{tabular}
\end{table}

\begin{figure}[H]
\centering
\includegraphics[width=\textwidth]{paper_figures/fig16_edge}
\caption{On-device cost on Samsung A17 5G. (a) Measured file sizes:
unified student ONNX against the teacher files present on disk (pest
teacher file absent and excluded). (b) Per-inference latency by provider
mode (10 warmup $+$ 100 measured, batch 1): dots are means with p95
upper whiskers and open-circle p50 markers. CPU and NNAPI coincide
($\pm$2.5\%) with bit-identical predictions, indicating CPU fallback on
this SoC. Single device ($n{=}100$ per mode; p95 approximate).}
\label{fig:edge}
\end{figure}

% ============================================================
% 13. DISCUSSION
% ============================================================

\section{Discussion}
\label{sec:discussion}

The held-out test results establish the system as a working unified
model: the proposed configuration scores 100.00\% leaf, 99.56\% pest,
and 95.69\% fruit accuracy with counting MAE 3.53 ($r{\approx}0.85$).
The story of the experiments, however, is not that every component
helps. The no-distillation ablation matches or beats the proposed
configuration on all five test metrics (fruit 99.14\%, MAE 3.08), and
the frozen backbone ties it on fruit test (95.69\%) at lower counting
error. Under the evaluated conditions -- teachers saturated near
ceiling tasks, distillation strengths fixed at $\alpha{=}\beta{=}0.5$
without sweeping, one seed per configuration -- heterogeneous
distillation does not improve the student, and LoRA adaptation is
neutral (its ${\sim}3$ pp fruit gap sits at the empirical run-to-run
noise floor estimated from the two no-LoRA configurations). The
positive ablation result is unambiguous: static weighting collapses
fruit test accuracy to 70.26\%, validating adaptive weighting as the
essential training component.

The learned task precisions also demonstrate that the optimization
process assigns different relative weights to the heterogeneous
objectives (fruit $\approx$ 18.8, leaf $\approx$ 7.0, pest
$\approx$ 5.7, yield $\approx$ 0.18 at the proposed best epoch).
These values should be interpreted as optimization parameters rather
than direct measures of task importance.

The ablation experiments therefore determine the contribution of each
training component rather than confirming prior expectations. The No-KD
experiment (\texttt{--alpha 1.0 --beta 0.0}) refutes the assumption
that teacher supervision helps here, Static Weights
(\texttt{--disable-homo}) confirms adaptive balancing is required,
No-LORA (\texttt{--no-lora}) shows adaptation is neutral within noise,
and No-Warm-up (\texttt{--freeze-ep 0}) shows staged optimization is
neutral. On-device, the unified model runs fully offline with
provider-identical numerics and no NNAPI acceleration on the test SoC
(Table~\ref{tab:edge}, Fig.~\ref{fig:edge}) -- efficiency, not
accuracy, is the contribution that survives contact with measurement.

% ============================================================
% 14. LIMITATIONS
% ============================================================

\section{Limitations and Future Work}
\label{sec:limitations}

Several experiments remain necessary before the study can be considered
complete.

First, the quantitative results for ablation configurations A--D must
be obtained from their corresponding final runs; configuration B has
only 103 logged epochs and must be completed.

Second, class-wise precision, recall, F1-score, and confusion matrices
are reported from held-out test predictions
(Tables~\ref{tab:leaf_classwise}--\ref{tab:fruit_classwise},
Figs.~\ref{fig:leaf_cm}--\ref{fig:fruit_cm}). Two caveats apply: no
Healthy fruit ROI is present in validation or test, so Healthy-fruit
generalization is untested; and the proposed configuration was
evaluated from the epoch-200 checkpoint (only
\texttt{latest\_checkpoint.pth} was present on the evaluation
machine) rather than the best-validation epoch 186.

Third, the fruit-counting branch is evaluated with
predicted-versus-ground-truth plots and correlation statistics
(Table~\ref{tab:counting}, Fig.~\ref{fig:counting}). The residual
limitation is the systematic undercounting bias ($-3.1$), which calls
for threshold calibration on dense clusters rather than further
ranking analysis.

Fourth, the edge evaluation is single-device (Samsung A17 5G):
a pilot on Samsung J6+ (Snapdragon 425, 4\,GB RAM, 2018) showed
${\sim}60$\,s per inference and was excluded from the comparison as
impractical for repeated benchmarking; ultra-low-end deployment
requires quantization and is deferred to future work. The latency tail
($n{=}100$ per mode) is approximate, battery readings are coarse drain
only with no wall-power claim, the latency subset is train-source (all
accuracy claims rest on the val accuracy set and the held-out college
test), and the pest phone images come from the training pool (reported
as laptop-agreement sanity while primary pest accuracy remains the
college held-out test).

Future work can investigate:

\begin{itemize}

\item post-training quantization,

\item quantization-aware training,

\item hardware-aware optimization,

\item additional field datasets,

\item robustness to illumination and occlusion,

\item domain adaptation,

\item temporal fruit tracking,

\item larger-scale orchard evaluation, and

\item comparison with additional single-task and multi-task baselines.

\end{itemize}

% ============================================================
% 15. CONCLUSION
% ============================================================

\section{Conclusion}
\label{sec:conclusion}

This paper presented Agri-CNX-Edge, a unified multi-task framework for
agricultural disease, pest, fruit classification, localization, and
counting.

The architecture combines a shared DINOv3-ConvNeXt-Tiny backbone with
LoRA adapters, hierarchical feature routing ($f1\!\rightarrow\!$leaf,
$f2\!\rightarrow\!$pest, $P3(f3{+}f4)\!\rightarrow\!$detection,
ROI-Align($P3$)$\!\rightarrow\!$fruit), lightweight MBConv
classification heads, a decoupled YOLO detection head with per-apple
ROI-Align classification, heterogeneous categorical and box-level
knowledge distillation, and adaptive homoscedastic task weighting.

On held-out test, the LoRA-adapted configuration reports 100.00\%
leaf accuracy, 99.56\% pest accuracy, and 95.69\% fruit accuracy. The
detection branch reports a Count MAE of 3.53 and Count RMSE of 6.13
($r{=}0.85$, $R^2{=}0.57$) on 400 yield images.

Against the frozen configuration (fruit test 95.69\%, MAE 3.43) and
the no-distillation ablation (fruit test 99.14\%, MAE 3.08), the
proposed configuration ties or trails on fruit and counting: the
measured contribution of this work is the unified efficiency system
(120.7\,MB, offline, provider-identical on-device numerics) and the
validation of adaptive weighting (static weighting collapses fruit
test accuracy to 70.26\%), alongside a documented negative result on
distillation under saturated-teacher conditions.

Class-wise test analysis, counting correlation statistics,
qualitative localization, and on-device efficiency are measured
(Tables~\ref{tab:leaf_classwise}--\ref{tab:edge},
Figs.~\ref{fig:leaf_cm}--\ref{fig:edge}). The remaining experimental
work is confirmatory: completing the static-weighting run past its
103 logged epochs, repeating configurations over seeds, and extending
on-device measurement to quantized builds.

% ============================================================
% REFERENCES
% ============================================================

\begin{thebibliography}{99}

\bibitem{kendall2018}
A. Kendall, Y. Gal, and R. Cipolla,
``Multi-Task Learning Using Uncertainty to Weigh Losses for Scene
Geometry and Semantics,''
in \textit{Proceedings of the IEEE/CVF Conference on Computer Vision
and Pattern Recognition (CVPR)}, 2018, pp. 7482--7491.

\bibitem{chen2018}
Z. Chen, V. Badrinarayanan, C.-Y. Lee, and A. Rabinovich,
``GradNorm: Gradient Normalization for Adaptive Loss Balancing in
Deep Multitask Networks,''
in \textit{Proceedings of the International Conference on Machine
Learning (ICML)}, 2018, pp. 794--803.

\bibitem{zhou2019}
X. Zhou, D. Wang, and P. Kr{\"a}henb{\"u}hl,
``Objects as Points,''
arXiv preprint arXiv:1904.07850, 2019.

\bibitem{caron2021}
M. Caron et al.,
``Emerging Properties in Self-Supervised Vision Transformers,''
in \textit{Proceedings of the IEEE/CVF International Conference on
Computer Vision (ICCV)}, 2021, pp. 9650--9660.

\bibitem{oquab2023}
M. Oquab et al.,
``DINOv2: Learning Robust Visual Features without Supervision,''
arXiv preprint arXiv:2304.07193, 2023.

\bibitem{liu2022}
Z. Liu et al.,
``A ConvNet for the 2020s,''
in \textit{Proceedings of the IEEE/CVF Conference on Computer Vision
and Pattern Recognition (CVPR)}, 2022, pp. 11976--11986.

\bibitem{howard2019}
A. Howard et al.,
``Searching for MobileNetV3,''
in \textit{Proceedings of the IEEE/CVF International Conference on
Computer Vision (ICCV)}, 2019, pp. 1314--1324.

\bibitem{tan2019}
M. Tan and Q. V. Le,
``EfficientNet: Rethinking Model Scaling for Convolutional Neural
Networks,''
in \textit{Proceedings of the International Conference on Machine
Learning (ICML)}, 2019, pp. 6105--6114.

\bibitem{hinton2015}
G. Hinton, O. Vinyals, and J. Dean,
``Distilling the Knowledge in a Neural Network,''
in \textit{NIPS Deep Learning and Representation Learning Workshop},
2015.

\bibitem{khan2022}
S. I. Khan, S. M. Quadri, S. Banday, and J. L. Shah,
``Deep diagnosis: A real-time apple leaf disease detection system
based on deep learning,''
\textit{Computers and Electronics in Agriculture},
vol. 198, 2022, Art. no. 107093.

\bibitem{gao2022}
F. Gao et al.,
``A novel apple fruit detection and counting methodology based on
deep learning and trunk tracking in modern orchard,''
\textit{Computers and Electronics in Agriculture},
vol. 197, 2022, Art. no. 107000.

\bibitem{janarthan2022}
S. Janarthan et al.,
``P2OP--Plant Pathology on Palms: A deep learning-based mobile
solution for in-field plant disease detection,''
\textit{Computers and Electronics in Agriculture},
vol. 202, 2022, Art. no. 107371.

\bibitem{huang2023}
Q. Huang et al.,
``Knowledge Distillation Facilitates the Lightweight and Efficient
Plant Diseases Detection Model,''
\textit{Plant Phenomics},
vol. 5, 2023, Art. no. 0062.

\bibitem{dhaka2021}
V. S. Dhaka et al.,
``A Survey of Deep Convolutional Neural Networks Applied for
Prediction of Plant Leaf Diseases,''
\textit{Sensors}, 2021.

\bibitem{liu2021}
J. Liu and X. Wang,
``Plant diseases and pests detection based on deep learning:
a review,''
\textit{Plant Methods},
vol. 17, 2021, Art. no. 22.

\bibitem{liu2021cagrad}
B. Liu, X. Liu, X. Jin, P. Stone, and Q. Liu,
``Conflict-Averse Gradient Descent for Multi-task Learning,''
in \textit{Advances in Neural Information Processing Systems
(NeurIPS)}, 2021.

\end{thebibliography}

\end{document}